\item El silbato de un tren ($f = 400\text{ Hz}$) suena más alto o más bajo en frecuencia dependiendo de si se aproxima o se aleja.
\begin{enumerate}
    \item Demuestre que la diferencia en frecuencia del silbato, si se acerca o se aleja es
$$ \Delta f = \frac{2(u/v)}{1 - u^2/v^2} f $$
donde $u$ es la rapidez del tren y $v$ es la rapidez del sonido.
    \item Calcule esta diferencia para un tren moviéndose con una rapidez de 130 km/h. Considere que la rapidez del sonido en aire es 340 m/s.
\end{enumerate}